\item \textbf{Absence of Licensed Legal Practitioner Evaluation (Most Critical Limitation for Deployment):} Due to geographical and institutional resource constraints, this study did not include evaluation by licensed legal practitioners. All faithfulness and safety assessments rely on automated LLM-as-a-judge scoring (Section~\ref{main-subsec:eval-metrics}) and continuous lexical metrics (ROUGE-L) evaluated against researcher-constructed reference answers cross-referenced against public legal reference sources (Cornell LII, Oyez, CRS Constitution Annotated). While the adversarial critic persona exhibits a modest, statistically significant correlation with reference lexical grounding (Pearson $r = \JudgePearsonR, p = \JudgePearsonP$), automated scoring cannot replicate practitioner judgment on subtle constitutional doctrine. Specifically, practitioners would assess legal correctness, temporal consistency, citation validity, and hallucination behavior beyond continuous metrics. Furthermore, lacking an independent labeled validation set for the NLI Smart Grader (Section~\ref{main-subsubsec:terra-nli-grader}), the false-non-entailment rate remains empirically unquantified beyond ablation evidence. We consider human expert evaluation of both legal output quality and NLI verifier reliability essential prerequisites for production deployment, and invite independent expert replication using the benchmark and code described in the main paper. The author spot-check described in Section~\ref{main-subsec:obscure-generalization} is a single non-expert review of a stratified sample and does not substitute for licensed practitioner review. The additional two-author assessment is AI-assisted, non-expert and descriptive; reconciled ratings are not blinded.

\item \textbf{Metric Disentanglement and Refusal-Volume Masking:} Evaluating complex RAG pipelines solely via aggregate faithfulness induces severe metric distortion: systems with high epistemic refusal volume can receive high faithfulness marks for non-committal answers on unanswerable queries, masking generative weaknesses. In the historical five-system dual-judge evaluation, TERRA has the highest overall ($\TerraOverallFaithfulnessMean$ vs.\ $\HippoRAGOverallFaithfulnessMean$; intervals overlap) and answer-only faithfulness ($\TerraAnswerOnlyFaithfulnessMean$ vs.\ $\HippoRAGAnswerOnlyFaithfulnessMean$) with fewer flagged refusals ($\TerraFlaggedRefusals$ vs.\ $\HippoFlaggedRefusals$), but answer-only scores are conditional on different answered subsets. On multi-hop reasoning TERRA leads HippoRAG by $\CatBTerraVsHippoRAGDelta$ ($q=\CatBTerraVsHippoRAGAdjQ$). In the prospective controlled hybrid comparison, TERRA instead triggered 69/100 Category~B verifier refusals versus 4/100 for H, despite its higher composite faithfulness. Neither score certifies useful-answer coverage or legal correctness. Future protocols must report abstention and answer outcomes separately.

\item \textbf{Evaluation rubric resolution and ceiling compression:} Coarse scales can obscure doctrinal differences. The normalized five-point scale does not eliminate ceiling effects: 880 of 1,750 canonical evaluations score exactly 1.0 (50.3\%). Automated rubric scores are not legal ground truth.

\item \textbf{Jurisdictional and Temporal Corpus Boundaries:} The pre-built Event Evolution Graph is anchored in \CorpusTotalOpinions\ authentic substantive U.S. Supreme Court opinions spanning \CorpusTemporalSpan, encompassing seven doctrinal domains of federal jurisprudence (Table~\ref{main-tab:corpus_stats}). Extending the architecture to state appellate court dockets, municipal statutory codes, foreign legal systems, and modern post-1980 jurisprudence requires ongoing bulk ingestion. Furthermore, while Section~\ref{main-subsec:temporal-update-stability} empirically demonstrates graph traversal stability across temporal snapshot expansions ($T \to T+N$), live vector index insertion and streaming graph maintenance remain unaddressed engineering extensions.

\item \label{item:hosted-api-determinism}\textbf{Hosted API Output Variability at $T=0.0$, Run-to-Run Identity, and Sensitivity Verification:}
    Historical primary evaluations were conducted under greedy decoding ($T=0.0$) using \texttt{databricks-meta-llama-3-3-70b-instruct} and \texttt{databricks-gemma-3-12b}. While setting $T=0.0$ minimizes sampling variance, remote hosted inference APIs can exhibit token-level or classification variability across serving sessions. Across benchmark evaluations, query \texttt{D066} exhibited an empirical classification shift (EASY $\to$ HARD), consistent with numerical or serving-level variability in hosted inference; the precise backend cause cannot be isolated from available logs. Comparing two 14-September TERRA executions across the 350-query benchmark suite, exact string-identical generated answers occurred on \TerraIdenticalAnswerFraction\ of queries overall (206/350), with pronounced category divergence: \TerraIdenticalAnswerCatA\ on Category~A (95/100) and \TerraIdenticalAnswerCatC\ on Category~C (62/75), versus \TerraIdenticalAnswerCatD\ on Category~D (39/75) and \TerraIdenticalAnswerCatB\ on Category~B (10/100). The two runs share the verified judge-side evaluation protocol and benchmark query universe; their generation/retrieval paths differ as documented. Both runs were evaluated under the standardized dual-judge protocol (\texttt{databricks-qwen3-next-80b-a3b-instruct} and \texttt{databricks-gpt-oss-120b}) using the frozen RAG prompt template (SHA256: \texttt{6d18fc14\dots}) and normalized Likert rubric. Sensitivity verification on an 80B model (\texttt{databricks-qwen3-next-80b-a3b-instruct}) across a separate 15-query subset (10 Cat~B, 5 Cat~D; distinct from the firewall stress test in Section~\ref{main-subsec:benchmark-design}) confirmed invariant relative rankings: TERRA achieved $1.000$ Category~B Faithfulness (vs. $0.900$ Flat RAG, $0.800$ Direct LLM) and $100\%$ Category~D safety rejection (vs. $40\%$ Flat RAG, $0\%$ Direct LLM). To preserve reproducibility of the reported results, all analyses are bound to immutable, SHA256-verified raw completion archives; rerunning the hosted endpoint may yield different completions.

\item \textbf{Parametric knowledge and temporal containment:} Direct LLM scored higher than TERRA on Category D continuous faithfulness ($\Delta=-0.200$, $p=3.8\times10^{-5}$, $q=1.27\times10^{-4}$; main FDR table, Rank 6) and Category B ($\CatBTerraVsDirectLLMDelta$, $q=\CatBTerraVsDirectLLMAdjQ$). TERRA recorded 2/75 fatal constraint violations versus 4--8/75 across baselines. Category D refusal and safe-handling comparisons form one eight-test Holm-adjusted secondary family; no safe-handling contrast survives that adjustment. Famous-case pretraining is a possible explanation, not an identified cause. TERRA's 73/75 safe outcomes comprise 48 refusals and 25 refutations in four-class scoring; system-flagged refusals are a different measure (46/75). D033 and D034 transposed the roles of Baker and Reynolds. Retrieval provenance does not purge post-cutoff knowledge from model weights or guarantee temporal correctness. Missing temporal judge values can reflect filtering or API errors; missingness alone does not establish the cause. On the obscure suite Direct LLM safe handling was \ObscDirectSafePct\ versus \ObscCrossDirectFamousSafe\ on the famous subset. Citation coverage and full-path recovery are retrieval/attribution measures, not proof of legal correctness. Adaptive invalidation-path traversal, weight-level approaches and expert checks remain future work.

\item \textbf{Fixed Graph Traversal Horizon and Multi-Hop Boundary Limits:} TERRA's online retrieval operates with a calibrated 2-hop BFS graph expansion radius to bound latency and token overhead. While 2-hop traversal captures direct citing and invalidating authorities across most evolutionary arcs, it can miss subsequent modifying precedents when citation chains exceed two steps. An empirical boundary emerged during our CaseFacts audit on \textit{Chimel v. California} (query \texttt{CF\_INC\_058}; Section~\ref{main-subsec:casefacts-generalization}), where the modifying precedent (\textit{United States v. Robinson} [1973]) sat beyond the 2-hop boundary; the generator cited \textit{Chimel}'s 1969 wingspan doctrine but found no modifying precedent in context. A limited retrieval horizon is a possible explanation for the gold-precedent omission rate during drafting on CLERC ($\ClercTerraMisapplicationTrackOne$; Section~\ref{main-subsec:clerc-generalization}), where incomplete context may encourage analogical reasoning from historical cases rather than fabrication. Developing adaptive traversal horizons along invalidation paths represents a high-priority architectural extension.

\item \textbf{Oracle Relation Labels, Exact-Title Seeding, and a Bundled Ablation:} In Full TERRA, curated \texttt{OVERRULES} edges are serialized with an explicit relation tag, dedicated heading, expanded excerpt budget, and prompt instruction; obscure benchmark pairs were drawn from these edges. The Full-TERRA advantage on that suite is therefore contingent on curated annotations, and the bundled ablation cannot isolate the tag alone. For named-case queries, seeding relied on exact caption matches (category-filtered vector retrieval returned no candidates), causing two queries to produce empty context and safe-rated refusals, one to omit the overruling case, and two to take the fast path; we evaluated TERRA as deployed and report sensitivity subsets excluding these items. The primary comparison contrasts complete systems, and Direct LLM receives no instruction block. Outcomes are judge-dependent in magnitude, the suite has thirty items, and judges were not fully blind to arm identity.

\item \label{item:falserefusal}\textbf{False Refusals on Answerable Queries (Categories A and B):}
In the historical canonical evaluation, the safety gate produces system-flagged refusals on answerable benchmark inquiries: \TerraCatAFalseRefusals/100 Category~A (Factual) and \TerraCatBFalseRefusals/100 Category~B (Evolutionary) queries receive system-level safety rejections (\texttt{safety\_rejected = True}). All \TerraCatBFalseRefusals\ Category~B false refusals originate from the NLI Smart Grader on the hard route; none occur on the fast route. Individual causes were not separately diagnosed. The prospective B refusal counts (T69/H4) are verifier decisions, not an independently adjudicated false-rejection rate.

\item \textbf{Latency measurement caveats:} Baselines ran sequentially on 5--6 September 2026; TERRA ran with four concurrent workers on 14 September. \TerraThrottledRows/350 TERRA queries incurred HTTP 429 back-off (\TerraThrottleSleepSeconds\,s total). EASY includes 110 queries (100 A and 10 D), gross/net median 5.465 s; HARD includes 240 (100 B, 75 C, 65 D), gross/net median 14.067 s. Category A/B medians (5.828/17.939 s) are not route medians. Retry sleep is subtracted per query before computing net percentiles. These are unlike harnesses; aggregate medians obscure routing and throttling, and comparisons are indicative only.

\item \label{item:samefamily}\textbf{Same-Family Generation, NLI Grading, and Category~D Verifier:}
Generation, NLI grading, and the Category~D verifier all deploy \texttt{databricks-\allowbreak meta-\allowbreak llama-\allowbreak 3-3-70b-\allowbreak instruct} (same family and endpoint). Faithfulness judges are \texttt{databricks-\allowbreak qwen3-next-\allowbreak 80b-a3b-\allowbreak instruct} and \texttt{databricks-\allowbreak gpt-oss-120b} (different model families). Category~D four-class labels (refusal, refutation, fatal error, hedge) may inherit shared Llama biases; cross-family verifier ablation has not been run.

\item \textbf{Unflagged Textual Refusals:}
System-level flagging (\texttt{safety\_rejected = True}) undercounts epistemic refusals. Adding textual refusals (asserting contextual absence without gate triggering) raises counts to \TerraEpistemicRefusals/350 for TERRA, \FlatEpistemicRefusals/350 for Flat Vector RAG, \GraphEpistemicRefusals/350 for GraphRAG-Adapted, and \HippoEpistemicRefusals/350 for HippoRAG-Adapted. Unflagged TERRA cases are D034 (refusal phrasing with fatal error, score 0.000), D061 (verified refutation, floored at \AmendmentMFloorScore\ under Amendment~M), and D068/D070 (fast-path textual refusals scoring 1.000). Aggregate Table~\ref{tab:sensitivity_raw_vs_clamped} sensitivity results for TERRA are invariant to this distinction ($\Delta = 0.000$).

\item \textbf{Mixed-date evaluation blocks:} Evaluation blocks differ in recorded dates, configurations and rubrics. The prospective hybrid comparison is separately frozen and does not recover the historical run identities. File modification dates are not execution timestamps. The obscure/famous records identify a recovered engine source and vector-index fingerprints, but exact historical code/graph/index linkage remains unavailable for the primary, pilot, CaseFacts and CLERC results. Temporal-probe payload hashes do not establish the generating engine or graph version. The pilot records identify a single Llama-3.3-70B judge, unlike canonical dual judging; task-specific scores are not interchangeable. Category B/D context profiles differ without establishing different execution dates. Requested generator endpoint parity does not establish backend weight invariance. See Table~\ref{tab:run_provenance}.

\item \textbf{Corpus Length-Filtering Bias and Concise Landmark Exclusion:}
Filtering by length ($\ge 4{,}000$ characters) purges noise: across \IssueTwentyTotalVolumesScanned\ volumes (\IssueTwentyTotalDecisionsScanned\ cases), \IssueTwentySubFiveHundredCharCount\ records are $<500$ characters (\IssueTwentySubFiveHundredCharPct, procedural orders/dockets). However, equating length with authority induces bias: it excludes \IssueTwentyIntermediateExcludedCount\ decisions ($500$--$3{,}999$ characters, \IssueTwentyIntermediateExcludedPct), with \IssueTwentyIntermediatePageRankGeNinetyNineFive\ in the top $0.5\%$ of CAP PageRank authority. Deterministic audit of 63 landmarks reveals \IssueTwentyLandmarksExcludedByLength\ (\IssueTwentyLandmarksExcludedByLengthPct) were purged, including procedural standards (\textit{Haines}, \IssueTwentyHainesCharCount\ chars, \IssueTwentyHainesCAPPRPercentile th percentile; \textit{Diebold}, \IssueTwentyDieboldCharCount\ chars, \IssueTwentyDieboldCAPPRPercentile th percentile), competency (\textit{Dusky}, \IssueTwentyDuskyCharCount\ chars, \IssueTwentyDuskyCAPPRPercentile th percentile), tax milestones (\textit{Lucas v. Earl}, \IssueTwentyLucasCharCount\ chars), and desegregation orders (\textit{Sipuel}, \IssueTwentySipuelCharCount\ chars; \textit{Gayle}, \IssueTwentyGayleCharCount\ chars). \textit{Bolling v. Sharpe} (\IssueTwentyBollingCharCount\ chars) met the threshold but was omitted in subsampling. Future graphs should replace static floors with citation retention.

\item \textbf{Citation Extraction Precision, Same-Year Chronology, and Provenance Uncertainty:}
Automated citation extraction via \texttt{eyecite} establishes citation existence rather than doctrinal endorsement. Our stratified audit ($N=\IssueTwentyOneAuditSampleN$, seed \IssueTwentyOneAuditSeed) confirms \IssueTwentyOneAuditCitationPresentPct\ of sampled edges contain reporter strings; direct precedential reliance is verified in \IssueTwentyOneAuditTargetConfirmedPct\ (95\% CI: \IssueTwentyOneAuditTargetConfirmedCI), with \IssueTwentyOneAuditQuotationPct\ in quotations. Furthermore, \IssueTwentyOneSameYearEdges\ edges (\IssueTwentyOneSameYearEdgesPct) connect same-year opinions, with \IssueTwentyOneIntraYearReversals\ intra-year date reversals from publication artifacts and zero cross-year forward violations. The legacy claim of ``$99.91\%$ chronological validity / $0.09\%$ noise'' in earlier working drafts was untraceable to repository artifacts, mandating this audit. Citation trails provide candidate authorities with \IssueTwentyOneAuditTargetConfirmedPct\ precision rather than certified doctrinal paths.

\item \textbf{Rule-Based Doctrinal Taxonomy, Lexical Polysemy, Truncation, and Residual Fallback Bias:}
The seven-domain classification is generated by a deterministic \IssueTwentyTwoTotalLexiconPhrases-phrase rule-based lexicon over truncated texts ($\le \IssueTwentyTwoTextTruncationChars$ characters), inducing four limitations: (1)~\textit{Lexical polysemy}: Keyword matching without contextual disambiguation risks misclassification when terms appear in procedural dicta or dissenting quotes. (2)~\textit{Truncation artifacts}: Truncating texts to \IssueTwentyTwoTextTruncationChars\ characters restricts matching to introductory syllabi, occasionally missing dispositive holdings. (3)~\textit{Residual fallback bias}: The \texttt{general\_constitutional\_statutory} category acts as an unkeyed fallback for \IssueTwentyTwoGeneralConstitutionalCorpus\ decisions (\IssueTwentyTwoGeneralConstitutionalCorpusPct) lacking substantive keywords, functioning as a catch-all with \IssueTwentyTwoGeneralConstitutionalBaselineAcc\ baseline precision despite \IssueTwentyTwoFallbackValidityPct\ validity as a residual class. (4)~\textit{Absence of human expert annotation}: Taxonomy validation relies on dual automated evaluators (IEA \IssueTwentyTwoIEARawAgreement, Cohen's $\kappa=\IssueTwentyTwoIEACohensKappa$) and a blinded automated adjudicator (\IssueTwentyTwoHorvitzThompsonAccuracy\ weighted accuracy, 95\% CI: \IssueTwentyTwoHorvitzThompsonCI), not licensed practitioner annotation; classifications are not certified human ground truth.

\item \textbf{A Priori Hyperparameter Selection and Sensitivity Analysis Scope:} Core retrieval and graph expansion parameters ($k=\IssueTwentyFourSeedK$, depth $d=\IssueTwentyFourBfsDepth$, hub cap $\IssueTwentyFourHubCap$, geometric threshold $\tau=\IssueTwentyFourDistThreshold$) were fixed prior to benchmark construction and evaluation based on topological network constraints and preliminary feasibility testing on an independent \IssueTwentyFourDevQueryCount-query development suite (\IssueTwentyFourDevDate). The development audit does not establish global optimality or exclude all contamination. Empirical sensitivity sweeps across \IssueTwentySixConfigCount\ configurations (Section~\ref{main-subsec:retrieval-quality}, Table~\ref{main-tab:retrieval_eval} Panel~C) profile performance across horizons ($k \in \{1,2,4\}$, $d \in \{1,2,3\}$, $C \in \{25,50,100\}$), describing the older stored-context configuration, not exact reproduction of canonical-v2 ($k=2, d=2, C=50$, \IssueTwentySixHistMacro\ macro recall, \IssueTwentySixHistFullPath\ full path recall) while quantifying the coverage-efficiency trade-off under bounded operational caps (\IssueTwentySixGoldRetentionPct\ landmark retention, \IssueTwentySixContextReductionPct\ context reduction). Canonical-v2 recovers 66/100 complete targets, older stored contexts 62/100 and the separate bounded engine 61/100; these are distinct records and configurations.

\item \textbf{Task-Corpus Domain Asymmetry and Private-Law Generalizability Scope:} A core methodological boundary is the intentional asymmetry between corpus composition and benchmark task scope. While the indexed repository reflects the complete 20,000-case docket spanning seven domains (\IssueTwentyFiveCorpusCombinedPrivateStatutoryPct\ private and general statutory law), the evaluated multi-hop reasoning benchmark (Category~B) evaluates constitutional and public-law precedent exclusively (\IssueTwentyFiveCatBConstitutionalPublicPct\ across five public-law domains; zero private-law chains). While the evaluation demonstrates TERRA's measured performance on the evaluated constitutional citation chains across a dense background network containing \CorpusPrivateLawOpinions\ private-law distractors, the system's performance on complex multi-hop statutory or patent reasoning remains unvalidated. Multi-hop evolutionary retrieval over complex commercial codes and specialized regulatory dockets constitutes an important open frontier for future domain-specific evaluation.
